% Authoritative LaTeX manuscript; Markdown is a generated compatibility projection.

% <!-- part-title: Modeling -->

% <!-- chapter-title: Structural Models: What Is Connected to What? -->

\chapter{Structural Models: What Is Connected to What?}

A structural model shows which parts of a system exist and how they relate to one another. We can ask, for example, which components depend on a service or where a document may be changed. To answer either question, we can set aside most of what happens inside the components.

\section{Entities, Relationships, and Graph Analysis}\label{heading:entities-relationships-and-graph-analysis}

We will examine each class of model in the same way. First, we construct a small model to answer a particular engineering question. We then examine the properties and analyses that this kind of model supports more generally. Finally, we turn to DocAble to see how the same ideas apply in a more complex system.

\subsection{Constructing a Structural Model}\label{heading:constructing-a-structural-model}

Consider a web application with an API, a cache, and a database. We want the web component to access the database only through the API, while allowing the cache to read from the database. The API provides a common place to apply access-control and consistency rules; a bypass could circumvent those rules.

% <!-- point: a-structural-model-selects-both-entities-and-relations-and-the-edge-label-carries-the-meaning | The architectural question selects components and their permitted direct dependencies. -->

\textbf{Question.} Can Web reach Database without passing through API? To answer, we represent the components and their dependencies in \autoref{canonical-structural}. We can leave out the cache's eviction policy and the database schema. Neither changes which components Web can reach in this model.

\begin{figure}[htbp]
\centering
\includegraphics[width=\linewidth]{book/latex/assets/canonical-structural-e9fe021e86.pdf}
\caption{\emph{Structural modeling---layered architecture.} This model states an intended architecture, rather than recording an execution. Every solid arrow denotes a direct dependency; the crossed arrow marks a forbidden bypass. Checking paths makes the obligation precise: every path from Web to Database must pass through API.}
\label{canonical-structural}
\end{figure}

\textbf{Semantics.} In \autoref{canonical-structural}, each solid arrow denotes a permitted direct dependency between two components. A path is a sequence of components connected by those arrows. The crossed arrow marks a forbidden direct dependency and is excluded when tracing paths. Web therefore has an indirect route to Database through API and Cache.

% <!-- point: once-structure-is-explicit-the-properties-are-graph-properties-and-the-analyses-graph-operations | Removing API and searching for a remaining Web-to-Database path tests whether the declared architecture permits a bypass. -->

\textbf{Analysis.} Our architectural requirement is stronger than prohibiting a direct Web-to-Database dependency: every path from Web to Database must pass through API. Remove API and the arrows touching it from a copy of the graph. Start at Web and follow the remaining arrows. In this example Web has no outgoing arrow, so Database is unreachable. Now imagine adding a permitted dependency from Web to Cache. Removing API would leave the path Web~\(\to\)~Cache~\(\to\)~Database, showing exactly how the revised architecture would permit a bypass.

\subsection{Properties and Analyses}\label{heading:structural-properties-and-analyses}

Structural models support several kinds of questions about how a system is organized. Depending on the relationships represented, we can investigate which components can reach one another, whether dependencies form cycles, how components are grouped, and whether relationships satisfy cardinality constraints.

\textbf{Reachability} asks whether a path connects one entity to another. \textbf{Separation} can require that no path connect two entities, or that they become disconnected when a designated boundary is removed. In \autoref{canonical-structural}, Web is allowed to reach Database, but every route must pass through API. These properties help engineers identify who can reach a protected resource and whether a proposed architecture permits a bypass. \textbf{Acyclicity} requires a dependency graph to contain no directed cycles. In a finite, acyclic build graph, the tasks can be ordered so that every prerequisite finishes before its dependent begins. If two compilation tasks each require the other's output, neither can go first. The dependencies must change before such an order is possible. \textbf{Containment} describes how entities are grouped, for example which modules belong to a service. \textbf{Cardinality} constrains how many relationships an entity may have, as when every service must have exactly one designated owner.

The corresponding analyses are graph operations. Follow dependency arrows forward to establish reachability or backward to identify which components a change might affect. Remove a boundary node and search for a remaining path to test mandatory traversal. A cycle check determines whether a prerequisite order exists, and counting the ownership links of each service identifies any with missing or multiple owners. Each analysis depends on the meaning of the relationships: the dependency graph in \autoref{canonical-structural} can answer the bypass question, but checking ownership requires ownership links.

For structural modeling, the procedure is to identify the entities and relationships relevant to the engineering question, define what those relationships mean, and analyze the resulting graph. We now turn to two structural models from DocAble: one describes where document mutation is permitted, and the other describes dependencies among remediation computations.

\section{DocAble: Architecture and Computation Dependencies}\label{heading:docable-architecture-and-computation-dependencies}

DocAble is a service that repairs accessibility problems in documents, such as missing alternative text, unmarked table headers, or incorrect reading order. The same operation may need a different implementation for a Word document and a PDF. Its structural models let engineers separate two questions: where format-specific changes belong, and how the repair computations depend on one another.

\subsection{Where may a document change?}\label{heading:where-may-a-document-change}

DocAble supports several document formats through libraries powerful enough to alter their underlying structures. Its architecture separates ordinary remediation logic from format-specific mutation through structured-document models. The Office formats use three models over a shared OpenXML foundation; PDF uses a separate model. The boundary is enforced differently across these implementations, but the architectural intent is common: document changes should pass through sanctioned mutation interfaces rather than arbitrary calls into format libraries.

\textbf{Document-mutation model · Structural}

\begin{itemize}
\tightlist
\item
  \textbf{Engineering question} --- Where is document mutation supposed to occur?
\end{itemize}

\begin{itemize}
\tightlist
\item
  \textbf{Model} --- a layered structural relation: remediation logic operates through a structured-document model, with a separate implementation for each format: DOCX, PPTX, XLSX, and PDF.
\end{itemize}

\begin{itemize}
\tightlist
\item
  \textbf{Property} --- the intended architecture places format-specific mutation behind one sanctioned seam; ordinary remediation logic is not intended to bypass it.
\end{itemize}

\begin{itemize}
\tightlist
\item
  \textbf{Quality attribute} --- architectural integrity and modifiability.
\end{itemize}

% <!-- point: a-structural-model-names-the-sanctioned-seam-for-document-mutation | The model draws one architectural relation — remediation reaches format-specific mutation through a structured-document seam, implemented per format — and leaves queues, retries, sessions, and repair mechanics out. -->

The necessary model is small. \autoref{document-mutation-model} draws it: remediation logic operates \emph{through} a structured-document model, which has a separate \emph{implementation} for each format---DOCX, PPTX, XLSX, PDF, and others. It says nothing about queues, retries, deployment, user sessions, or the mechanics of individual repairs. Those facts are real but irrelevant to the question; the reduction preserves only the architectural boundary that matters.

\begin{figure}[htbp]
\centering
\includegraphics[width=\linewidth]{book/latex/assets/document-mutation-model-450f92b9ce.pdf}
\caption{\emph{Document-mutation structural model.} The reduction preserves one architectural relation: ordinary remediation reaches format-specific mutation through a structured-document seam. Enforcement of that relation is a separate Alignment decision.}
\label{document-mutation-model}
\end{figure}

% <!-- point: a-model-precedes-the-constraint-that-protects-it | The model states the path requirement that ordinary mutation routes cross the structured-document seam; DocAble protects the boundary using format-specific static restrictions and runtime write-authority checks. | terms: soft-vs-hard-governance -->

The principal property is the mandatory-traversal constraint introduced in \autoref{canonical-structural}: every ordinary route to format-specific mutation must cross a sanctioned seam. In a sufficiently detailed dependency graph, we could remove that seam and ask whether remediation logic can still reach a raw mutation facility. Any remaining path would identify a bypass. DocAble instead uses format-specific static restrictions and runtime write-authority checks to protect these boundaries. \magevalue{chapter:alignment} examines those controls.

% <!-- point: a-structural-model-can-drift-from-the-system-it-describes | A structural model can become stale as implementation changes — the code moves and the drawing does not — so the boundary it shows quietly stops matching the system; that gap is model drift, and later sections keep the correspondence by deriving, generating, or tracing. -->

% <!-- index-def: model-drift -->

A structural model can become stale as the implementation changes. The code moves and the drawing does not, so the boundary it shows quietly stops matching the system. That gap between a model and the system it describes is \textbf{model drift}. Later sections make the correspondence explicit by deriving, generating, or tracing between model and implementation.

\subsection{What computations depend on one another?}\label{heading:what-computations-depend-on-one-another}

% <!-- point: a-second-structural-question-reaches-inside-the-pipeline | A second structural question reaches inside remediation: which computations make up the pipeline, and how do they depend on one another; the first reduction treats each registered pass as a node. -->

The document-mutation model identifies where changes belong. A second structural question concerns how remediation computations depend on one another. If we change one computation, which others might be affected? DocAble runs documents through registered remediation passes that detect conditions, produce information for later passes, decide whether work should proceed, or modify document state. A \emph{registered computation} is a named remediation pass that DocAble knows how to invoke, such as a pass that classifies document content or repairs alternative text. The registry identifies the pass and records selected properties of its role in the pipeline. The graph represents the pass as a node, leaving its internal statements, library calls, and loops out of the graph.

A repair pass may edit a document directly or produce a \textbf{typed patch}: a structured description of a proposed change, such as assigning alternative text to a figure. A separate component, the \textbf{applier}, interprets the patch and performs the change. The boundary where proposed changes become document mutations is the \textbf{apply seam}.

\textbf{Remediation-computation model · Structural}

\begin{itemize}
\tightlist
\item
  \textbf{Engineering question} --- What computations make up remediation, and how does information move among them?
\end{itemize}

\begin{itemize}
\tightlist
\item
  \textbf{Model} --- a static directed graph whose nodes are registered computations and whose typed edges represent declared composition among them.
\end{itemize}

\begin{itemize}
\tightlist
\item
  \textbf{Property} --- Declared data and control dependencies are distinguishable and consistent with pipeline ordering. Each computation declares its mutation behavior: read-only, direct editing, typed patch production, or a combination of patch production and direct editing. Patch producers also identify where their patches are applied.
\end{itemize}

\begin{itemize}
\tightlist
\item
  \textbf{Quality attribute} --- analyzability, modifiability, and architectural integrity.
\end{itemize}

% <!-- point: the-computation-model-distinguishes-data-flow-from-control-composition | DocAble's computation model represents registered computations as nodes and distinguishes a data-flow edge (information consumed to produce a result) from a control edge (information consulted only to decide whether to run). -->

DocAble\textquotesingle s remediation-computation model distinguishes two relationships between registered computations. A data-flow edge means one computation produces information another consumes in producing its result. A control edge means one computation produces information another consults only to decide whether it should run. \autoref{from-nodes-to-composition} gives an example of each kind of \emph{typed composition}: the classifier determines which repair passes run, while the repair passes produce patches that the patch applier consumes. The first relationship governs execution and the second carries data. Recording the edge type preserves that distinction. These distinctions make impact questions straightforward to investigate: the document-signals classifier has fifteen immediate control-gate dependents whose execution decisions may change with its output, while the intra-page vision-fusion pass has two immediate data-flow consumers. An engineer or agent can inspect these edges directly or traverse them to investigate indirect dependencies. Neither requires a dedicated analysis subsystem.

\begin{figure}[htbp]
\centering
\includegraphics[width=\linewidth]{book/latex/assets/from-nodes-to-composition-1d52db233f.pdf}
\caption{\emph{From computations to composition.} This illustrative fragment distinguishes data flow (solid arrows) from control dependencies (dashed arrows). Ellipses identify read-only computations, diamonds typed patch producers, and rounded boxes direct document edits. The apply seam is the boundary where patches are applied; the final alt-text check reports findings without changing the document.}
\label{from-nodes-to-composition}
\end{figure}

% <!-- point: the-model-also-records-how-bounded-each-effect-is | The same model records each computation's mutation behavior — read-only, a typed patch for all or part of its edits, or a direct in-place edit — exposing distinctions the implementation contains but does not represent uniformly. -->

The model also declares each computation\textquotesingle s mutation behavior. Some computations only inspect the document, some produce typed patches, some produce patches for part of their work and edit directly for the rest, and most edit document state directly. Patch producers also declare the seam through which their patches are applied. These attributes let an engineer ask which computations can modify a document directly and which produce patches for an applier. They describe the intended behavior of each computation; separate checks establish how closely the implementation follows those declarations.

% <!-- point: the-static-graph-is-not-one-execution-of-it | The computation model is static per run, but other views attach observations to its stable computation identities — DocAble joins staged-run latency and cost to those identities without turning the structural graph into an execution log. -->

The computation model describes the declared computations and dependencies. Per-job invocations and retries belong to execution records. Other views can nevertheless attach observations to its stable computation identities. DocAble joins staged-run telemetry to those identities, so measured latency and cost can be queried alongside the structural model without turning the structural graph into an execution log. The per-session mutation record remains separate: it records what happened to one artifact, not the execution history of the computation graph.

% <!-- point: the-identity-spine-supports-added-views-without-enlarging-the-schema | Additional engineering views attach by computation identity rather than enlarging the structural schema; declared properties state intent, measured properties report observation, so an unmeasured value is unknown, not zero. -->

The same identities can support additional engineering views without enlarging the structural schema each time. Separate providers can record complexity, bounded execution, latency, test coverage, or checker trust against the same computation identity. Declared properties state engineering intent. Measured properties report observations and may have incomplete coverage. An unmeasured latency is therefore unknown, not zero. Joining these views lets engineers ask questions about combinations of properties while keeping each representation focused on its own question.

% <!-- point: the-contrast-between-the-two-structural-models | The document-mutation model locates architectural bypasses; the computation model supports change-impact analysis through data and control dependencies. -->

The two structural models answer different questions about the same software. To investigate a possible bypass of the document abstraction, examine the architectural boundary and the paths that could circumvent it. To investigate a change to a remediation pass, follow its declared data and control dependencies and inspect its mutation behavior. Neither analysis requires a complete model of DocAble.

The two DocAble examples illustrate why structural models must preserve the relationships relevant to their questions. Architectural paths, data dependencies, and control dependencies preserve different facts even when each is drawn with boxes and arrows. The next section adds time: what states a thing occupies, and what sequences of changes the model admits.
